\documentclass[a4paper,11pt]{article}
\usepackage{exam-style-341}
\renewcommand{\examtitle}{\thispagestyle{empty}\begin{center}\vspace*{1.5cm}{\fontsize{20pt}{28pt}\selectfont\bfseries\color{ctp-text}341 农业综合知识三}\\[0.3cm]{\fontsize{14pt}{18pt}\selectfont\bfseries\color{ctp-mauve}综合诊断卷（一）参考答案}\end{center}\vspace{0.8cm}\hrule\vspace{0.8cm}\setcounter{page}{1}}
\fancyhead[L]{\fontsize{9pt}{10pt}\selectfont\color{ctp-subtext0}341 农业综合知识三 · 参考答案一}\fancyhead[R]{\fontsize{9pt}{10pt}\selectfont\color{ctp-subtext0}C 语言 · 网络 · 数据库}
\begin{document}\examtitle
\parttitle{第一部分 C 语言}
\q 空字符（\code{\\0}）。\quad \q 指针。\quad \q \code{fclose}。\quad \q 0；1。\quad \q 终止（基例）。\par
\q 值传递复制实参，形参修改不影响实参；指针传递地址值，通过解引用可修改原对象。\quad \q 数组是连续同类型元素集合，数组名在多数表达式中转换为首元素指针，但数组本身不是可修改的指针变量。\quad \q 文本方式按字符并可能转换换行，二进制方式按字节原样读写。\par
\q 输出：\code{1 5 4 9}。\quad \q 输出：\code{24}。\quad \q 交换后主函数中的\code{x,y}互换，因为传入了它们的地址。\par
\q 参考实现：先求最大值，再在小于最大值的元素中取最大者；\code{a==NULL}、\code{result==NULL}、\code{n<2} 或不存在不同第二大值时返回 0。
\begin{lstlisting}[language={[custom]C}]
int second_max(const int a[], int n, int *result) {
    int i, max, second, found = 0;
    if (!a || !result || n < 2) return 0;
    max = a[0];
    for (i=1; i<n; ++i) if (a[i] > max) max = a[i];
    for (i=0; i<n; ++i)
        if (a[i] < max && (!found || a[i] > second))
            second = a[i], found = 1;
    if (!found) return 0; *result = second; return 1;
}
\end{lstlisting}
\parttitle{第二部分 计算机网络}
\q 7；32；IP；可靠；MAC。\quad \q 分层把复杂通信划分为模块，各层通过接口协作；发送端逐层加首部封装，接收端反向剥离解封装。\quad \q TCP 可靠有序、面向连接，适合网页和文件；UDP 无连接、开销小、时延低，适合 DNS 和实时音视频。\quad \q IP 标识网络层主机接口，MAC 标识链路层接口，端口号标识主机中的进程。\quad \q HTTPS/HTTP 产生数据，TCP 形成报文段，IP 形成数据报，以太网形成帧；交换机依据 MAC 表，路由器依据目标 IP 和路由表；HTTPS 通过 TLS 提供加密、认证和完整性。
\parttitle{第三部分 数据库技术与应用}
\q DBMS；元组；\code{HAVING}；持久性；2NF。\quad \q DB 是数据集合，DBMS 是管理数据库的软件，DBS 是数据库、DBMS、应用和人员等组成的系统。\quad \q 外模式/概念模式映射保障逻辑独立性，概念模式/内模式映射保障物理独立性。\par
\q \code{SELECT Sname FROM S WHERE Sdept='信息系';}\quad \code{SELECT Sno,AVG(Grade) FROM SC GROUP BY Sno;}\quad\code{PRIMARY KEY(Sno,Cno), FOREIGN KEY(Sno) REFERENCES S(Sno), FOREIGN KEY(Cno) REFERENCES C(Cno)}。\quad 候选键为(学号,课程号)，存在重复、插入、删除和修改异常；分解为\code{学生(学号,姓名)}、\code{课程(课程号,课程名,任课教师)}、\code{选课(学号,课程号,成绩)}。
\end{document}
